% ------------------------------------------------
% RESUMEN
% ------------------------------------------------

\begin{abstract}
El recocido simulado es una de las metaheurísticas más usadas para el
problema del agente viajero (TSP) y, al mismo tiempo, una cadena de
Markov cuyo comportamiento puede analizarse con rigor. En este trabajo
se construye la cadena de Metropolis sobre el espacio de rutas con la
vecindad 2-opt y se demuestra que es irreducible, aperiódica y
reversible respecto de la distribución de Boltzmann. Se estudia además
cómo la temperatura controla a la vez la concentración de esa
distribución en las rutas óptimas y la velocidad de convergencia, medida
por la brecha espectral. Con ello, el recocido aparece como una cadena
no homogénea cuya convergencia caracteriza el teorema de Hajek. También
se analiza una variante sobre la formulación QUBO del TSP y se demuestra
que la penalización de las restricciones crea barreras de energía que
crecen linealmente con el parámetro \(M\).

Los experimentos, reproducibles en un notebook de Google Colab,
confirman la teoría. Con cinco ciudades, las frecuencias empíricas
convergen a la distribución exacta y la pendiente de \(-\ln\gamma(T)\)
frente a \(1/T\) es \(68.2\), a un \(3\,\%\) de la profundidad crítica
exacta \(h^\ast=66\). El recocido sobre rutas encuentra el óptimo en
las \(40\) corridas sobre cuatro instancias de TSPLIB, hasta
\texttt{berlin52}. En cambio, sobre el QUBO el desempeño depende de
\(M\) en forma de U y, en \texttt{burma14}, no baja del \(29\,\%\) por
encima del óptimo. La conclusión es que, para una misma cadena de
Metropolis, la elección del espacio de estados determina la geometría
del paisaje y, con ella, la eficacia del algoritmo.

\medskip
\noindent\textbf{Palabras clave:} cadenas de Markov, algoritmo de
Metropolis, recocido simulado, brecha espectral, problema del agente
viajero, QUBO.
\end{abstract}
